% ===== enrich_ : Module & taxonomy validation (Phase 4) =====
\subsection{Biological annotation of the leading harmonics and the compartment layer}
\label{sec:enrich}

We tested whether the leading omics harmonics correspond to coherent biological
programmes, and whether the five-compartment / seven-class layer is anchored to
ground-truth annotation. For each dataset we aggregated the per-slice loading
magnitude $\langle|\varphi_{n,i}|\rangle_s$ of the three leading harmonics and ran
hypergeometric over-representation analysis (ORA) on the top-40 loading genes against
nine hand-curated gene sets (circadian clock, cell cycle, OXPHOS/energy,
immune/inflammation, sterol/lipid biosynthesis, redox/antioxidant, ECM/adhesion,
IEG/signalling, xenobiotic metabolism; sets defined explicitly in the analysis code),
using all mapped features as the background and Benjamini--Hochberg FDR correction.
Enrichment was fully offline; no live annotation service was queried.

\paragraph{The leading harmonic $\varphi_1$ is an expression-magnitude artefact, not a
programme.} On \emph{neither} dataset did the top-loading genes of $\varphi_1$ enrich
any gene set (minimum adjusted $p = 1.00$ on GSE10072; $0.77$ on GSE11923). Across all
three leading harmonics the circadian dataset (GSE11923) produced no significant term
(minimum adjusted $p = 0.77$). The reason is mechanical: because the omics connectome
matrix is amplitude-weighted ($H_{ij}=c_{ij}\,\psi_i\psi_j^\ast$), $|\varphi_{1,i}|$
tracks mean log-expression per feature (Spearman $\rho = 0.43$ on GSE10072, $\rho =
0.42$ on GSE11923; both $p<10^{-17}$). The dominant mode therefore reports which genes
are highly expressed, not which genes co-regulate.

\paragraph{Core-clock genes do \emph{not} load preferentially on $\lambda_1$.} The
central prediction --- that in circadian data the clock module dominates the leading
harmonic --- is refuted. On GSE11923 core-clock genes load \emph{below} background on
$\varphi_1$ (one-sided Mann--Whitney for clock $>$ background $p = 0.996$, rank-biserial
$= -0.23$); the same holds on GSE10072 ($p > 0.999$, rank-biserial $= -0.40$)
(Fig.~\ref{fig:enrich_clock}). Clock genes carry moderate, not extreme, expression and
so sit in the body of the $|\varphi_1|$ distribution.

\paragraph{Sub-leading modes recover annotation-supported biology --- on the cancer
dataset only.} Once the amplitude-dominated $\varphi_1$ is set aside, coherent
programmes appear in the sub-leading harmonics of GSE10072: $\varphi_2$ enriches the
redox/antioxidant set ($k=9$, adjusted $p = 2.9\times10^{-3}$) and the circadian-clock
set (adjusted $p = 1.4\times10^{-2}$), and $\varphi_3$ enriches IEG/signalling ($k=9$,
adjusted $p = 5.1\times10^{-3}$) (Fig.~\ref{fig:enrich_ora}). Consistently, four of the
five compartment marker sets load preferentially on $\varphi_2$ in GSE10072 (Clock,
Redox, Signalling, Biosynthesis all Mann--Whitney $p<0.01$; Energy $p=0.10$). No
sub-leading harmonic reached significance on the circadian dataset, and no compartment
marker set loaded preferentially there.

\paragraph{The compartment layer is descriptive, not a validated classifier.}
The evidence does not support a mechanistic reading of the leading mode or of a
seven spectral state classes. The leading harmonic is an expression-magnitude summary; the
clock-loading prediction does not hold on the circadian dataset; and the
only annotation-supported structure is confined to sub-leading modes and
present in one of the two datasets. We therefore treat the Compartment Coupling
Matrix and the spectral state classes as a descriptive decomposition rather than a
mechanistic classifier: the sub-leading harmonics of the cancer dataset provide an
annotation-consistent decomposition, but the leading mode and the five compartments
should not be read as a ground-truth-anchored biological classifier.

\begin{figure}[t]
  \centering
  \includegraphics[width=\linewidth]{enrich_harmonic_enrichment.png}
  \caption{Hypergeometric over-representation of nine curated gene sets among the top-40
  loading genes of the three leading harmonics ($\varphi_1$ blue, $\varphi_2$ orange,
  $\varphi_3$ green), for the cancer (GSE10072, left) and circadian (GSE11923, right)
  datasets. Marker size is proportional to the overlap count; black-outlined markers are
  significant at BH-adjusted $p<0.05$ (dashed line). The leading harmonic $\varphi_1$
  enriches nothing on either dataset; significant terms appear only in the sub-leading
  modes of the cancer dataset.}
  \label{fig:enrich_ora}
\end{figure}

\begin{figure}[t]
  \centering
  \includegraphics[width=0.85\linewidth]{enrich_clock_loading.png}
  \caption{Leading-harmonic loading $|\varphi_{1,i}|$ of core-clock genes (red) versus
  all other features (grey); black bars are group medians. In both the circadian
  (GSE11923, left) and cancer (GSE10072, right) datasets clock genes load at or below
  background, refuting the prediction that the clock module dominates $\lambda_1$
  (one-sided Mann--Whitney $p$ and rank-biserial correlation annotated).}
  \label{fig:enrich_clock}
\end{figure}
